\BourbakiExercisesSection

\begin{enumerate}
\def\labelenumi{\arabic{enumi}.}
\tightlist
\item
  \begin{enumerate}
  \def\labelenumii{(\alph{enumii})}
  \tightlist
  \item
    Let E be a right A-module. On the additive group \(E^{r}\) (with \(r \geqslant 1\) ) an external law of composition is defined with \(\mathbf{M}_{r}(\mathbf{A})\) as set of operators, denoting by x.P, for every element \(x = (x_{i})_{1 \leqslant i \leqslant r}\) of \(E^{r}\) and every square matrix \(P = (\alpha_{ij}) \in \mathbf{M}_{r}(\mathbf{A})\) , the element \(y = (y_{i})\) of \(E^{r}\) such that
  \end{enumerate}
\end{enumerate}

\[
y _ {i} = \sum_ {j = 1} ^ {r} x _ {j} \alpha_ {j i} \quad (1 \leqslant i \leqslant r).
\]

This external law defines with the additive law on \(E^{r}\) a right \(\mathbf{M}_{r}(\mathbf{A})\) -module structure on \(E^{r}\) ; by restricting the ring of operators of the ring of scalar matrices \(I.\alpha\) , the product A-module structure on \(E^{r}\) is recovered. Show that for the A-module E to have a generating system of r elements, it is necessary and sufficient that the \(\mathbf{M}_{r}(\mathbf{A})\) -module \(E^{r}\) be monogenous.

\begin{enumerate}
\def\labelenumi{(\alph{enumi})}
\setcounter{enumi}{1}
\tightlist
\item
  Let \((x_{i})_{1\leqslant i\leqslant n}\) be a generating system of E and \((y_{j})_{1\leqslant j\leqslant m}\) a family of elements of E (resp. a generating system of E). Let \(z, z', z''\) be three elements of \(\mathbf{E}^{m + n}\) such that \(z_{i} = x_{i}\) for \(1\leqslant i\leqslant n\) , \(z_{n + j} = 0\) for \(1\leqslant j\leqslant m\) , \(z_{i}' = x_{i}\) for
\end{enumerate}

\(1 \leqslant i \leqslant n, z_{n+j}^{\prime} = y_{j}\) for \(1 \leqslant j \leqslant n, z_{i}^{\prime \prime} = 0\) for \(1 \leqslant i \leqslant n, z_{n+j}^{\prime \prime} = y_{j}\) for \(1 \leqslant j \leqslant n\) . Show that there exist two invertible matrices \(P, Q\) of \(\mathbf{M}_{n+m}(\mathbf{A})\) such that \(z' = z.P\) (resp. \(z'' = z.Q\) ).

\begin{enumerate}
\def\labelenumi{(\alph{enumi})}
\setcounter{enumi}{2}
\tightlist
\item
  If A is commutative, \(u_{P}: x \mapsto x \cdot P\) is an endomorphism of the A-module \(E^{r}\) and the mapping \(P \mapsto u_{P}\) is a homomorphism of the ring \(\mathbf{M}_{r}(\mathbf{A})\) into the ring \(\operatorname{End}_{\mathbf{A}}(\mathbf{E}^{r})\) . If E is a faithful A-module this homomorphism is injective.
\end{enumerate}

\begin{enumerate}
\def\labelenumi{\arabic{enumi}.}
\setcounter{enumi}{1}
\tightlist
\item
  \begin{enumerate}
  \def\labelenumii{(\alph{enumii})}
  \tightlist
  \item
    Let \(X\) be a square matrix over a ring \(\mathbf{A}\) , which can be written as an upper triangular matrix \((X_{ij})\) of matrices \((1 \leqslant i \leqslant p, 1 \leqslant j \leqslant p)\) . Show that if each of the square matrices \(X_{ii}\) \((1 \leqslant i \leqslant p)\) is invertible, so is \(X\) and \(X^{-1}\) can be written as an upper triangular matrix \((Y_{ij})\) \((1 \leqslant i \leqslant p, 1 \leqslant j \leqslant p)\) corresponding to the same partition of the indexing set. When \(\mathbf{A}\) is a field, prove that this sufficient condition for \(X\) to be invertible is also necessary.
  \end{enumerate}
\end{enumerate}

\begin{enumerate}
\def\labelenumi{(\alph{enumi})}
\setcounter{enumi}{1}
\tightlist
\item
  Give an example of a ring A and a matrix \(\begin{pmatrix} a & 1 \\ 0 & b \end{pmatrix}\) (with \(a \in \mathbf{A}, b \in \mathbf{A}\) ) which is invertible without either \(a\) or \(b\) being invertible in A and whose inverse \(\begin{pmatrix} a' & b' \\ c' & d' \end{pmatrix}\) is such that \(b' \neq 0\) and \(c' \neq 0\) (take A to be the endomorphism ring of an infinite-dimensional vector space).
\end{enumerate}

\begin{enumerate}
\def\labelenumi{\arabic{enumi}.}
\setcounter{enumi}{2}
\tightlist
\item
  Let A be the quotient ring Z/30Z. Show that in the matrix
\end{enumerate}

\[
\left( \begin{array}{c c c} 1 & 1 & - 1 \\ 0 & 2 & 3 \end{array} \right)
\]

over the ring A, the two rows are linearly independent but any two columns are linearly dependent.

\begin{enumerate}
\def\labelenumi{\arabic{enumi}.}
\setcounter{enumi}{3}
\tightlist
\item
  Let A be a ring, C its centre, B the ring of matrices \(\mathbf{M}_n(\mathbf{A}), \Delta\) the additive subgroup of A generated by the elements \(\alpha\beta - \beta\alpha\) for \(\alpha \in \mathbf{A}, \beta \in \mathbf{A}\) , and D the additive subgroup of B generated by the matrices \(XY - YX\) for \(X \in \mathbf{B}, Y \in \mathbf{B}\) ; \(\Delta\) and D are C-modules. For every matrix \(X = (\xi_{ij}) \in \mathbf{B}\) , let \(\theta(X)\) be the element \(\sum_{i=1}^{n} \bar{\xi}_{ij}\) of A/ \(\Delta\) , where, for all \(\alpha \in \mathbf{A}\) , \(\bar{\alpha}\) denotes the class of \(\alpha\) mod. \(\Delta\) . Show that \(\theta\) is a surjective homomorphism of B onto A/ \(\Delta\) , whose kernel is equal to D, so that B/D is isomorphic to A/ \(\Delta\) as a C-module. (Observe that D contains the matrices \(\alpha E_{ij}\) and the matrices \(\alpha(E_{ii} - E_{jj})\) for \(\alpha \in \mathbf{A}\) and \(i \neq j\) .)
\end{enumerate}

¶ 5. Let A be a ring, L, M, N any three indexing sets \(U = (a_{\lambda \mu})\) a matrix of \(\mathbf{A}^{\mathbf{L} \times \mathbf{M}}\) and \(V = (b_{\mu \nu})\) a matrix of \(\mathbf{A}^{\mathbf{M} \times \mathbf{N}}\) ; if, for every ordered pair \((\lambda, \nu) \in \mathbf{L} \times \mathbf{N}\) , the family \((a_{\lambda \mu} b_{\mu \nu})\) ( \(\mu \in \mathbf{M}\) ) has finite support, the element \(c_{\lambda \nu} = \sum_{\mu} a_{\lambda \mu} b_{\mu \nu}\) is defined and it is also said that the matrix \((c_{\lambda \nu})\) is the product \(UV\) of \(U\) by \(V\) . When the products \(UV'\) and \(UV''\) are defined, so is \(U(V' + V'')\)

and \(UV' + UV'' = U(V' + V'')\) ; is the converse true? Give an example of three infinite matrices, U, V, W such that the products UV, VW, \(U(VW)\) and \((UV)W\) are defined but \(U(VW) \neq (UV)W\) (take U to be a matrix with one row all of whose elements are equal to 1, W its transpose and V a matrix all of whose elements are 0, 1 or -1 and which has only a finite number of elements \(\neq 0\) in each row and each column. Determine by induction the elements of V such that UV = 0 but VW has a single element \(\neq 0\) .)

\begin{enumerate}
\def\labelenumi{\arabic{enumi}.}
\setcounter{enumi}{5}
\item
  Let X be a matrix with m rows and n columns over a field K; show that the rank \(\operatorname{rg}(X)\) is equal to the greatest of the ranks of the submatrices of X with an equal number of rows and columns. (Let \(\operatorname{rg}(X) = r\) ; if \(a_{1}, \ldots, a_{r}\) are r columns of X forming a free system in \(K_{d}^{m}\) and a basis of \(K_{d}^{m}\) is formed with these r vectors and m - r vectors of the canonical basis, show that the components of \(a_{1}, \ldots, a_{r}\) over the r other vectors of the canonical basis form a matrix of rank r.)
\item
  Let X be a matrix with m rows and n columns over a field K; if r is the rank of X, show that the rank of a submatrix with m rows and s columns, obtained by suppressing n - s columns of X, is \(\geqslant r + s - n\) .
\item
  Let \(X = (\alpha_{ij})\) be a matrix with \(m\) rows and \(n\) columns over a field \(K\) . For \(X\) to be of rank 1, it is necessary and sufficient that there exist in \(K\) a family \((\lambda_i)_{1 \leqslant i \leqslant m}\) of \(m\) elements not all zero and a family \((\mu_j)_{1 \leqslant j \leqslant n}\) of \(n\) elements not all zero such that \(\alpha_{ij} = \lambda_i \mu_j\) for every ordered pair of indices.
\item
  Let X, Y be two matrices with m rows and n columns over a field K; if there exist two square matrices \(P, P_{1}\) of order m and two square matrices \(Q, Q_{1}\) of order n such that Y = PXQ and \(X = P_{1}XQ_{1}\) , show that X and Y are equivalent.
\item
  Let \(X, X', Y, Y'\) be four square matrices of order \(n\) over a ring \(A\) , such that \(X\) is invertible. For there to exist two invertible square matrices \(P, Q\) of order \(n\) such that \(X' = PXQ\) and \(Y' = PYQ\) , it is necessary and sufficient that \(X'\) be invertible and that the matrices \(YX^{-1}\) and \(Y'X'^{-1}\) be similar.
\item
  Let E be a right vector space over a field K, of dimension \(\geqslant1\) , and H a hyperplane of E. Every endomorphism u of E leaving invariant each of the elements of H gives, on passing to the quotients, an endomorphism of the 1-dimensional quotient space E/H, an endomorphism which is therefore of the form \(\dot{x}\mapsto\dot{x}\mu(\dot{x})\) , where \(\mu(\dot{x})\in\mathrm{K}\) is such that \(\mu(\dot{x}\lambda)=\lambda^{-1}\mu(\dot{x})\lambda\) for \(\lambda\in K^{*}\) . An automorphism of E leaving invariant the elements of H is called a transvection of hyperplane H if the corresponding automorphism of E/H is the identity and a dilatation of hyperplane H otherwise; when u is a dilatation, the set of elements \(\mu(\dot{x})\) for \(\dot{x}\in E/H\) which is a class of conjugate elements in the multiplicative group \(K^{*}\) , is called the class of the dilatation u.
\end{enumerate}

\begin{enumerate}
\def\labelenumi{(\alph{enumi})}
\item
  Show that for every dilatation there exists one and only one supplementary line of H, invariant under the dilatation.
\item
  Let \(\phi\) be a linear form on \(\mathbf{E}\) such that \(\mathbf{H} = \overline{\phi}(0)\) ; show that for every transvection \(u\) of hyperplane \(\mathbf{H}\) there exists a unique vector \(a \in \mathbf{H}\) such that \(u(x) = x + a\phi(x)\) for all \(x \in \mathbf{E}\) . Let \(\Gamma(\mathbf{E}, \mathbf{H})\) be the subgroup of \(\mathbf{GL}(\mathbf{E})\) consisting of the automorphisms leaving invariant each element of \(\mathbf{H}\) ; show that the transvections of hyperplane \(\mathbf{H}\) form a normal commutative subgroup \(\Theta(\mathbf{E}, \mathbf{H})\) of \(\Gamma(\mathbf{E}, \mathbf{H})\) isomorphic to the additive group \(\mathbf{H}\) ; the quotient group \(\Gamma(\mathbf{E}, \mathbf{H}) / \Theta(\mathbf{E}, \mathbf{H})\) is isomorphic to the multiplicative group \(\mathbf{K}^*\) . For \(\Gamma(\mathbf{E}, \mathbf{H})\) to be commutative, it is necessary and sufficient that one of the two following conditions be satisfied: ( \(\alpha\) ) \(\mathbf{K}\) is a field with two elements; ( \(\beta\) ) \(\mathbf{K}\) is commutative and \(\dim \mathbf{E} = 1\) .
\item
  Suppose that E is finite-dimensional; show that for every transvection u there exists a basis of E such that the matrix of u with respect to this basis has all its diagonal elements equal to 1 and at most one other element \(\neq0\) .
\item
  Show that the centralizer (I, § 5, no. 3) of the group \(\Theta(E, H)\) in the group \(\mathbf{GL}(E)\) is the composition \(Z(E)\,\Theta(E, H) = \Theta(E, H)Z(E)\) of \(\Theta(E, H)\) and the centre \(Z(E)\) of \(\mathbf{GL}(E)\) (cf.~§ 1, Exercise 26). The only automorphisms belonging to this centralizer and leaving invariant at least one element \(\neq0\) of \(E\) are the transvections of the group \(\Theta(E, H)\) . If \(K\) has at least 3 elements, the centralizer of \(\Gamma(E, H)\) in \(\mathbf{GL}(E)\) is equal to \(Z(E)\) .
\item
  Show that the normalizers (I, § 5, no. 3) of \(\Theta(\mathbf{E},\mathbf{H})\) and \(\Gamma(\mathbf{E},\mathbf{H})\) in \(\mathbf{GL}(\mathbf{E})\) are both equal to the subgroup consisting of the automorphisms leaving H invariant.
\end{enumerate}

¶ 12. Let E be a right vector space of dimension ≥1 over a field K. Let F(E) denote the normal subgroup of GL(E) consisting of the automorphisms u such that the kernel of 1\_E - u is of finite codimension (hence F(E) = GL(E) if E is finite-dimensional). Let SL(E) denote the normal subgroup of GL(E) generated by all the transvections (Exercise 11); it is contained in F(E).

\begin{enumerate}
\def\labelenumi{(\alph{enumi})}
\item
  If \(\dim E \geqslant 2\) , show that for every ordered pair of non-zero vectors x, y of E, there exists a transvection or product of two transvections, which maps x to y (in other words \(\mathbf{SL}(\mathbf{E})\) operates transitively on \(E = \{0\}\) ).
\item
  Let V, W be two hyperplanes of E, \(\dot{x}_{0} = x_{0} + V\) a class mod. V distinct from V and \(\dot{y}_{0} = y_{0} + W\) a class mod. W distinct from W. Show that if \(\dim E \geqslant 2\) , there exists a transvection or product of two transvections, which maps V to W and \(\dot{x}_{0}\) to \(\dot{y}_{0}\) (consider first the case where V and W are distinct).
\item
  If \(\dim \mathbf{E} \geqslant 2\) , show that any two transvections distinct from the identity are conjugate in the group \(\mathbf{F}(\mathbf{E})\) .
\item
  If \(\dim \mathbf{E} \geqslant 3\) , show that any two transvections distinct from the identity are conjugate in the group \(\mathbf{SL}(\mathbf{E})\) (reduce it with the aid of (b) to the case where the hyperplanes of the two transvections are identical, then use (a)).
\item
  Suppose that \(\dim E = 2\) . For any two transvections to be conjugate in \(\mathbf{SL}(\mathbf{E})\) , it is necessary and sufficient that the subgroup \(\mathbf{Q}\) of \(\mathbf{K}^*\) generated by the squares of elements of \(\mathbf{K}^*\) be identical with \(\mathbf{K}^*\) . (If \(u\) is a transvection distinct from the identity, \(a\) a vector of \(\mathbf{E}\) which is not invariant under \(u\) and \(b = u(a) - a\) , show that, for every transvection \(u'\) conjugate to \(u\) in \(\mathbf{SL}(\mathbf{E})\) , \(u'(a) - a = a\lambda + b\mu\) , with \(\mu = 0\) or \(\mu \in \mathbf{Q}\) ; for this use the fact that in a group \(\mathbf{G}\) every product \(sts^{-1}t\) is a product of squares.)
\item
  Suppose that \(\dim E \geqslant 2\) . For two dilatations \(u, u'\) to be such that there exists a \(v \in \mathbf{SL}(\mathbf{E})\) such that \(u' = vuv^{-1}\) , it is necessary and sufficient that the classes (Exercise 11) of \(u\) and \(u'\) be the same (use \((b)\) ). Deduce that if the class of a dilatation is contained in the commutator group of \(\mathbf{K}^*\) , this dilatation belongs to the group \(\mathbf{SL}(\mathbf{E})\) (observe that \((vuv^{-1})u^{-1} = v(uv^{-1}u^{-1}))\) .
\end{enumerate}

¶ 13. Let E be a right vector space of dimension ≥2 and H₀ a hyperplane of E.

\begin{enumerate}
\def\labelenumi{(\alph{enumi})}
\item
  Show that every automorphism \(u\) belonging to \(\mathbf{F}(\mathbf{E})\) (Exercise 12) is the product of an automorphism of \(\mathbf{SL}(\mathbf{E})\) and possibly a dilatation of the hyperplane \(\mathbf{H}_0\) (proceed by induction on the codimension of the kernel of \(1_{\mathbf{E}} - u\) , using Exercises 12(a) and 12(f)).
\item
  Show that \(\mathbf{SL}(\mathbf{E})\) contains the commutator group of \(\mathrm{F(E)}\) and is identical with this group except when \(\mathbf{E}\) is a space of dimension 2 over the field with 2 elements. (To show that \(\mathbf{SL}(\mathbf{E})\) contains the commutator group of \(\mathrm{F(E)}\) , use (a) and Exercise \(12(f)\) . To see that \(\mathbf{SL}(\mathbf{E})\) is contained in the commutator group of \(\mathrm{F(E)}\) , except in the case indicated, show that for every homomorphism of \(\mathrm{F(E)}\) into a commutative group, the image of a transvection is the identity element, using Exercises \(11(b)\) and \(12(c)\) .)
\item
  Show that \(\mathbf{SL}(\mathbf{E})\) is equal to its commutator group except when \(\dim \mathbf{E} = 2\) and \(\mathbf{K}\) has 2 or 3 elements. (If \(\dim \mathbf{E} \geqslant 3\) , note that every transvection \(u\) can be written as \(vwv^{-1}w^{-1}\) , where \(v\) is a transvection and \(w \in \mathbf{SL}(\mathbf{E})\) , using Exercise 12(d). If \(\dim \mathbf{E} = 2\) , note first that every automorphism \(u\) whose matrix with respect to a basis of \(\mathbf{E}\) is of the form \(\begin{pmatrix} \lambda & 0 \\ 0 & \lambda^{-1} \end{pmatrix}\) is a product of transvections, arguing as in (a); then consider the commutator \(uvu^{-1}v^{-1}\) , where \(v\) is the transvection whose matrix with respect to the same basis is \(\begin{pmatrix} 1 & 1 \\ 0 & 1 \end{pmatrix}\) .)
\end{enumerate}

¶ 14. Let E be a vector space of dimension ≥2 over a field K.

\begin{enumerate}
\def\labelenumi{(\alph{enumi})}
\item
  Show that the projective group PGL(E) is canonically isomorphic to the quotient of the linear group GL(E) by its centre (isomorphic to the multiplicative group of the centre of K, cf.~§ 1, Exercise 26).
\item
  Let \(\mathbf{PSL}(\mathbf{E})\) denote the canonical image of the group \(\mathbf{SL}(\mathbf{E})\) in \(\mathbf{PGL}(\mathbf{E})\) ; it is a normal subgroup of \(\mathbf{PGL}(\mathbf{E})\) , containing the commutator group of \(\mathbf{PGL}(\mathbf{E})\) when \(\mathbf{E}\) is finite-dimensional. Show that \(\mathbf{PSL}(\mathbf{E})\) is a doubly transitive (cf.~I, § 5, Exercise 14) group of permutations of \(\mathbf{P}(\mathbf{E})\) .
\item
  Let \(a \neq 0\) be an element of \(\mathbf{E}\) and let \(e = \pi(a) \in \mathbf{P}(\mathbf{E})\) (in the notation of §9, no. 6). Show that the transvections of the form \(x \mapsto x + a\phi(x)\) (where \(\phi \in \mathbf{E}^*\) is such that \(\phi(a) = 0\) ) constitute a commutative subgroup of \(\mathbf{SL}(\mathbf{E})\) ; if \(\Gamma_e\) is the image of this subgroup in \(\mathbf{PSL}(\mathbf{E})\) , \(\Gamma_e\) is a normal subgroup of the subgroup \(\Phi_e\) of \(\mathbf{PSL}(\mathbf{E})\) leaving \(e\) invariant. Show also that the union of the subgroups conjugate to \(\Gamma_e\) in \(\mathbf{PSL}(\mathbf{E})\) generates \(\mathbf{PSL}(\mathbf{E})\) .
\item
  Let \(\Sigma\) be a primitive group (I, § 5, Exercise 13) of permutations of a set F which is equal to its commutator group. Suppose that there exists an element \(c \in \mathbf{F}\) such that the subgroup \(\Phi_c\) of \(\Sigma\) leaving \(c\) invariant contains a commutative normal subgroup \(\Gamma_c\) such that the union of the subgroups conjugate to \(\Gamma_c\) in \(\Sigma\) generates \(\Sigma\) . Show that under these conditions \(\Sigma\) is simple. (Let \(\Delta\) be a normal subgroup of \(\Sigma\) distinct from the identity; using I, § 5, Exercise 17, show that \(\Sigma = \Delta\) . \(\Phi_c = \Delta\) . \(\Gamma_c\) and, using the fact that \(\Gamma_c\) is commutative, deduce that every commutator in \(\Sigma\) is contained in \(\Delta\) .)
\item
  Deduce from (c), (d) and Exercise 13(c) that the group \(\mathbf{PSL}(\mathbf{E})\) is simple except when \(\dim \mathbf{E} = 2\) and \(\mathbf{K}\) has 2 or 3 elements.
\end{enumerate}

\(^{*}(f)\) If \(\dim E = n + 1\) and K is a field with q elements, show that the order of PGL(E) is

\[
(q ^ {n + 1} - 1) (q ^ {n + 1} - q) \dots (q ^ {n + 1} - q ^ {n - 1}) q ^ {n}
\]

(use (a), Exercise 9(a) of § 9, and the fact that K is necessarily commutative (V, § 11, Exercise 14 and VIII, § 11, no. 1, Theorem 1)).*

\begin{enumerate}
\def\labelenumi{(\alph{enumi})}
\setcounter{enumi}{6}
\item
  Show that if \(\dim \mathbf{E} = 2\) and \(\mathbf{K}\) is a field with 2 (resp. 3) elements, \(\mathbf{PSL}(\mathbf{E})\) is isomorphic to the symmetric group \(\mathfrak{S}_3\) (resp. the alternating group \(\mathfrak{A}_4\) ). (Considering \(\mathbf{PSL}(\mathbf{E})\) as a group of permutations of \(\mathbf{P}(\mathbf{E})\) , note that it contains every transposition (resp. every cycle ( \(a b c\) ) and use \(f\) ).)
\item
  Show that, except in the cases considered in (g), every normal subgroup of \(\mathbf{SL}(\mathbf{E})\) distinct from \(\mathbf{SL}(\mathbf{E})\) is contained in the centre of \(\mathbf{SL}(\mathbf{E})\) (use (d) and the fact that every transvection is a commutator in \(\mathbf{SL}(\mathbf{E})\) (Exercise 13(c))).
\end{enumerate}

¶ 15. Let A be a ring; for every integer \(n \geqslant 1\) , \(\mathbf{A}^{n}\) is canonically identified with the submodule of \(\mathbf{A}^{(\mathbf{N})}\) consisting of the elements whose coordinates of index \(\geqslant n\) are 0 and \(\mathbf{A}_{n}^{\prime}\) denotes the complementary submodule consisting of the elements whose coordinates of index \(<n\) are 0. Every endomorphism \(u \in \mathbf{GL}_{n}(\mathbf{A})\) is identified with the automorphism of \(\mathbf{A}^{(\mathbf{N})}\) whose restriction to \(\mathbf{A}^{n}\) is \(u\) and whose restriction to \(\mathbf{A}_{n}^{\prime}\) is the identity, so that \(\mathbf{GL}_{n}(\mathbf{A})\) is identified with a subgroup of \(\mathbf{GL}(\mathbf{A}^{(\mathbf{N})})\) ; let F denote the subgroup of \(\mathbf{GL}(\mathbf{A}^{(\mathbf{N})})\) the union of the \(\mathbf{GL}_{n}(\mathbf{A})\) , in other words the subgroup leaving invariant the elements of the canonical basis of \(\mathbf{A}^{(\mathbf{N})}\) except for a finite number of them (cf.~Exercise 12).

Let \(\mathbf{T}_n\) be the subgroup of \(\mathbf{GL}_n(\mathbf{A})\) generated by the \(\mathrm{B}_{ij}(\lambda)\) for \(0 \leqslant i < n\) , \(0 \leqslant j < n\) , \(i \neq j\) , \(\lambda \in \mathbf{A}\) (no. 13) and let \(\mathbf{T}\) be the subgroup of \(\mathbf{F}\) the union of the \(\mathbf{T}_n\) .

\begin{enumerate}
\def\labelenumi{(\alph{enumi})}
\item
  Show that for \(n \geqslant 3\) , \(T_n\) is equal to its commutator group.
\item
  Let \(a, b\) be two invertible elements of \(A\) . Show that
\end{enumerate}

\[
\left( \begin{array}{c c} a b & 0 \\ 0 & 1 \end{array} \right) = P \left( \begin{array}{c c} a & 0 \\ 0 & b \end{array} \right) Q \quad \text { and } \quad \left( \begin{array}{c c} a & 0 \\ 0 & b \end{array} \right) = P ^ {\prime} \left( \begin{array}{c c} b & 0 \\ 0 & a \end{array} \right) Q ^ {\prime}
\]

where \(P, Q, P', Q'\) belong to \(T_2\) . Deduce that the matrix \(\begin{pmatrix} aba^{-1}b^{-1} & 0 \\ 0 & 1\end{pmatrix}\) belongs to \(T_2\) .

\begin{enumerate}
\def\labelenumi{(\alph{enumi})}
\setcounter{enumi}{2}
\tightlist
\item
  Deduce from (b) that the commutator group of \(\mathbf{GL}_{n}(\mathbf{A})\) is contained in \(T_{2n}\) (replace A by \(\mathbf{M}_{n}(\mathbf{A})\) in (b)). Conclude (using (a)) that the commutator group of F is equal to T.
\end{enumerate}

\begin{enumerate}
\def\labelenumi{\arabic{enumi}.}
\setcounter{enumi}{15}
\tightlist
\item
  \begin{enumerate}
  \def\labelenumii{(\alph{enumii})}
  \tightlist
  \item
    Let \(U\) be a square matrix of order \(n\) over the field \(\mathbf{Q}\) all of whose non-diagonal elements are equal to the same rational number \(r > 0\) and whose diagonal \((d_1, \ldots, d_n)\) is such that \(d_i \geqslant r\) for all \(i\) and \(d_i > r\) except perhaps for one value of \(i\) . Show that \(U\) is invertible (if \(x = (x_i)_{1 \leqslant i \leqslant n}\) is a solution of the equation \(U.x = 0\) , show that the \(x_i\) all have the same sign and deduce that they are all zero).
  \end{enumerate}
\end{enumerate}

\begin{enumerate}
\def\labelenumi{(\alph{enumi})}
\setcounter{enumi}{1}
\tightlist
\item
  Let E be a finite set with m elements denoted by \(a_{i}\) ( \(1 \leqslant i \leqslant m\) ) and \((\mathbf{A}_{j})_{1 \leqslant j \leqslant n}\) a family of n subsets of E, distinct from one another; suppose that \(\text{Card}(\mathbf{A}_{i} \cap \mathbf{A}_{j}) = r \geqslant 1\) for every ordered pair \((i, j)\) of distinct elements of \([1, n]\) . Show that necessarily \(n \leqslant m\) . (Let \(V = (c_{ij})\) be the matrix of type \((m, n)\) such that \(c_{ij} = 1\) when \(a_{i} \in A_{j}\) , \(c_{ij} = 0\) otherwise. Apply (a) to the matrix \(^{t}V\) . V of order n.)
\end{enumerate}

¶ 17. Let K be a field (commutative or otherwise), E an n-dimensional right vector space over K and u an endomorphism of E.

\begin{enumerate}
\def\labelenumi{(\alph{enumi})}
\item
  Show that if \(n > 1\) and \(u\) is not a central homothety, there exists a basis of \(E\) such that the matrix of \(u\) with respect to this basis is of the form \((\alpha_{ij})\) with \(\alpha_{ij} = 0\) for \(j > i + 1\) and \(\sum_{i=1}^{n-1} \alpha_{i,i+1} = 1\) (argue by induction on \(n\) ).
\item
  Show that if \(n > 2\) , there exists a basis of \(\mathbf{E}\) such that the matrix of \(u\) with respect to this basis is of the form \((\alpha_{ij})\) with \(\alpha_{ij} = 0\) for \(j > i + 1\) and
\end{enumerate}

\(\sum_{i=1}^{n-1} \alpha_{i, i+1} = 0\) . (Argue by induction on \(n\) using \((a)\) , noting that it amounts to the same to prove the proposition for \(u\) or for \(u + \gamma 1_{E}\) , where \(\gamma\) belongs to the centre of \(K\) , and studying separately the case where \(u\) is of rank 1 or 2.)

\begin{enumerate}
\def\labelenumi{(\alph{enumi})}
\setcounter{enumi}{2}
\item
  Let \(A = (\alpha_{ij})\) be a matrix with the properties described in (b). Also let \(S\) be the matrix of order \(n(\varepsilon_{ij})\) such that \(\varepsilon_{ij} = 0\) except when \(i = j + 1\) , in which case \(\varepsilon_{j+1,j} = 1\) . Show that there exists a matrix \(B = (\beta_{ij})\) of order \(n\) such that \(A = \lambda E_{nn} + (BS - SB)\) with \(\lambda \in \mathbf{K}\) . (Take \(B\) such that \(\beta_{i1} = 0\) for \(1 \leqslant i \leqslant n\) and \(\beta_{ij} = 0\) for \(j > i + 2\) .)
\item
  Suppose that K is commutative. Deduce from (c) that every endomorphism u of E such that \(\operatorname{Tr}(u) = 0\) can be written as vw - wv, where v and w are endomorphisms of E.
\end{enumerate}

